%-%-%-%-%-%-%-%-%-%-%-%-%-%-%-%-%-%-%-%-%-%-%-%-%-%-%-%-%-%-%-%-%-%-%-%-%-%-%-%
%     ___                        _                           _  _   __        
%    /   \  ___    _ __    ___  | |_   _ __ ___    ___    __| |(_) / _| _   _ 
%   / /\ / / _ \  | '_ \  / _ \ | __| | '_ ` _ \  / _ \  / _` || || |_ | | | |
%  / /_// | (_) | | | | || (_) || |_  | | | | | || (_) || (_| || ||  _|| |_| |
% /___,'   \___/  |_| |_| \___/  \__| |_| |_| |_| \___/  \__,_||_||_|   \__, |
%                                                                       |___/ 
% 
%-%-%-%-%-%-%-%-%-%-%-%-%-%-%-%-%-%-%-%-%-%-%-%-%-%-%-%-%-%-%-%-%-%-%-%-%-%-%-%
\documentclass[a4paper,10pt,twoside]{article}
\usepackage{conf}
\begin{document}

%-%-%-%-%-%-%-%-%-%-%-%-%-%-%-%-%-%-%-%-%-%-%-%-%-%-%-%-%-%-%-%-%-%-%-%-%-%-%-%
%                     _  _   __        
%   /\/\    ___    __| |(_) / _| _   _ 
%  /    \  / _ \  / _` || || |_ | | | |
% / /\/\ \| (_) || (_| || ||  _|| |_| |
% \/    \/ \___/  \__,_||_||_|   \__, |
%                                |___/ 
% 
%-%-%-%-%-%-%-%-%-%-%-%-%-%-%-%-%-%-%-%-%-%-%-%-%-%-%-%-%-%-%-%-%-%-%-%-%-%-%-%

%-%-%-%-%-%-%-%-%-%-%-%-%-%-%-%-%-%-%-%-%-%-%-%-%-%-%-%-%-%-%-%-%-%-%-%-%-%-%-%
%  __   _                 _   
% / _\ | |_   __ _  _ __ | |_ 
% \ \  | __| / _` || '__|| __|
% _\ \ | |_ | (_| || |   | |_ 
% \__/  \__| \__,_||_|    \__|
% 
%-%-%-%-%-%-%-%-%-%-%-%-%-%-%-%-%-%-%-%-%-%-%-%-%-%-%-%-%-%-%-%-%-%-%-%-%-%-%-%

\conferenceinfo{62}{2026}

\udk{004.8}

\topic{Аппаратная реализация компактной и эффективной нейронной сети}

\information{Кривальцевич~Е.~А., студент гр.~555701}

\supervisor{Вашкевич~М.~И. -- докт.~техн.~наук}

%-%-%-%-%-%-%-%-%-%-%-%-%-%-%-%-%-%-%-%-%-%-%-%-%-%-%-%-%-%-%-%-%-%-%-%-%-%-%-%
%    _                         _           _    _               
%   /_\   _ __   _ __    ___  | |_   __ _ | |_ (_)  ___   _ __  
%  //_\\ | '_ \ | '_ \  / _ \ | __| / _` || __|| | / _ \ | '_ \ 
% /  _  \| | | || | | || (_) || |_ | (_| || |_ | || (_) || | | |
% \_/ \_/|_| |_||_| |_| \___/  \__| \__,_| \__||_| \___/ |_| |_|
% 
%-%-%-%-%-%-%-%-%-%-%-%-%-%-%-%-%-%-%-%-%-%-%-%-%-%-%-%-%-%-%-%-%-%-%-%-%-%-%-%
\annotation{В данной статье представлен пример шаблона оформления материала для участия в Научной конференции аспирантов, магистрантов и студентов БГУИР. Приведены примеры вставки изображений, таблиц и формул, а также даны рекомендации по оформлению ссылок и списка литературы.}

\keyword{СНТК, конференция, шаблон.}

\begin{multicols}{2} %НЕ УДАЛЯТЬ!

%-%-%-%-%-%-%-%-%-%-%-%-%-%-%-%-%-%-%-%-%-%-%-%-%-%-%-%-%-%-%-%-%-%-%-%-%-%-%-%
%  _____              _   
% /__   \  ___ __  __| |_ 
%   / /\/ / _ \\ \/ /| __|
%  / /   |  __/ >  < | |_ 
%  \/     \___|/_/\_\ \__|
% 
%-%-%-%-%-%-%-%-%-%-%-%-%-%-%-%-%-%-%-%-%-%-%-%-%-%-%-%-%-%-%-%-%-%-%-%-%-%-%-%
% Первый раздел
\section*{Введение}
В данной статье описаны основные элементы шаблона, которые могут быть использованы при подготовке материалов для Научной конференции аспирантов, магистрантов и студентов БГУИР. Шаблон обеспечивает единообразие оформления и упрощает процесс подготовки публикации.

Текст без формул и картинок помещается просто в исходный документ LaTeX без каких-либо специальных команд. Чтобы начать новый абзац~-- просто вставьте пустую строку. Заботиться о переносах не нужно. 

% Второй раздел
\section{Пример использования изображений}
Данный шаблон позволяет вставлять изображения как в ширину одного столбца, так и на всю ширину страницы.

Пример вставки изображения формата PNG в пределах одного столбца приведён ниже.
\image{
[width=7.8cm]{png/gr.png}
\caption{Пример изображения в пределах одного столбца}
}

Для вставки изображения, занимающего всю ширину страницы, можно использовать конструкцию, приведённую после основного текста статьи.

В тексте можно делать ссылки на источники, например: \ref{Vash2025}, на документацию по оформлению~\ref{STP}, а также на другие публикации~\ref{OtherRef}.

% Третий раздел
\section{Пример использования таблиц}
В данном разделе приведён пример оформления таблиц в стиле шаблона.

\myTable{Пример простой таблицы}{{\columnwidth}{|p{2cm}|X|X|X|} \hline
Параметр  & Значение~1 & Значение~2 & Значение~3 \\ \hline
А & 0.5 & 1 & 2.25 \\ \hline
Б & 0.0 & 3 & 9 \\ \hline
}
Ширина столбцов таблицы определяется автором. Важно, чтобы размер шрифта оставался неизменным и таблица не выходила за границы столбца.

\myTable{Пример объединения столбцов и строк}{
{\columnwidth}{|p{2.5cm}|X|X|X|} \hline
\multirow{2}{*}{Параметр} & \multicolumn{3}{|c|}{Значения} \\ \cline{2-4}
 & 1 & 2 &  3 \\ \hline
А & 0.5 & 1 & 2.25 \\ \hline
Б & 0.0 & 3 & 9 \\ \hline
\multirow{2}{*}{В} & 1.0 & \multicolumn{2}{|c|}{Совмещённые ячейки} \\ \cline{2-4}
 & 2.5 & 5.0 & 7.5 \\ \hline
}

% Четвертый раздел  
\section{Пример использования формул}
Текст в статье можно выделять \textbf{полужирным} или \textit{курсивным} начертанием для акцентирования внимания на ключевых элементах.

Формулы, вставленные в строку, оформляются с использованием конструкции \verb|\(\)|. Например:
\(\displaystyle a = (b + c)\cdot\mu\cdot\frac{\Theta}{\alpha^{3-i}}\).

Если формулу необходимо вынести на отдельную строку, используется конструкция \verb|\[\]|:
\[
a = \sum^9_{i=0}b\cdot \frac{k}{l}
\]

При необходимости можно добавить нумерацию уравнений — для этого используется окружение \verb|equation|.  

В результате формула будет отображена на отдельной строке и автоматически пронумерована:

\begin{equation}
E = m \cdot c^2
\label{eq:sum}
\end{equation}

Если требуется сослаться на формулу в тексте, нужно добавить метку с помощью команды \verb|\label{}| и ссылку на неё — \verb|\ref{}|.  

Результат в тексте будет выглядеть так:  
Как видно из уравнения~(\ref{eq:sum}), выражение задаёт зависимость между параметрами.

\section{Пример использования списков}
Списки позволяют компактно и наглядно структурировать информацию.

Ненумерованный список создаётся с помощью конструкции \verb|\begin{itemize}|:

Результат:
\begin{itemize}
  \item Первый элемент списка
  \item Второй элемент
  \item Третий элемент
\end{itemize}

Нумерованный список создаётся аналогично, но с окружением \verb|enumerate|:

Результат:
\begin{enumerate}
  \item Первый пункт
  \item Второй пункт
  \item Третий пункт
\end{enumerate}

\section{Пример вставки исходного кода}
Для вставки фрагментов программного кода используется пакет \texttt{listings}.  
Например, для языка Python:

\begin{lstlisting}[language=Python, caption={Пример Python-кода}]
def add(a, b):
    return a + b

print(add(3, 5))
\end{lstlisting}

%-%-%-%-%-%-%-%-%-%-%-%-%-%-%-%-%-%-%-%-%-%-%-%-%-%-%-%-%-%-%-%-%-%-%-%-%-%-%-%
%    ___                      _              _               
%   / __\  ___   _ __    ___ | | _   _  ___ (_)  ___   _ __  
%  / /    / _ \ | '_ \  / __|| || | | |/ __|| | / _ \ | '_ \ 
% / /___ | (_) || | | || (__ | || |_| |\__ \| || (_) || | | |
% \____/  \___/ |_| |_| \___||_| \__,_||___/|_| \___/ |_| |_|
% 
%-%-%-%-%-%-%-%-%-%-%-%-%-%-%-%-%-%-%-%-%-%-%-%-%-%-%-%-%-%-%-%-%-%-%-%-%-%-%-%
\section{Заключение}
В статье продемонстрированы основные возможности шаблона оформления материалов для конференции: добавление изображений, таблиц, формул и списка литературы. Соблюдение данных рекомендаций позволит обеспечить единый стиль публикаций и ускорить процесс подготовки тезисов.

%-%-%-%-%-%-%-%-%-%-%-%-%-%-%-%-%-%-%-%-%-%-%-%-%-%-%-%-%-%-%-%-%-%-%-%-%-%-%-%
%    ___  _  _      _  _                                    _            
%   / __\(_)| |__  | |(_)  ___    __ _  _ __   __ _  _ __  | |__   _   _ 
%  /__\//| || '_ \ | || | / _ \  / _` || '__| / _` || '_ \ | '_ \ | | | |
% / \/  \| || |_) || || || (_) || (_| || |   | (_| || |_) || | | || |_| |
% \_____/|_||_.__/ |_||_| \___/  \__, ||_|    \__,_|| .__/ |_| |_| \__, |
%                                |___/              |_|            |___/ 
% 
%-%-%-%-%-%-%-%-%-%-%-%-%-%-%-%-%-%-%-%-%-%-%-%-%-%-%-%-%-%-%-%-%-%-%-%-%-%-%-%

\section*{Список использованных источников}
Если при первой компиляции вместо ссылок на источники отображаются символы \texttt{??}, необходимо повторно скомпилировать проект. Это связано с тем, что файл \texttt{.aux}, в котором LaTeX хранит информацию о ссылках, обновляется не сразу.

\begin{ListReferences}
  \refitem{Vash2025}{Vashkevich M., Krivalcevich E. Compact and Efficient Neural Networks for Image Recognition Based on Learned 2D Separable Transform // 2025 27th International Conference on Digital Signal Processing and its Applications (DSPA). – IEEE, 2025. – С. 1–6.}
  \refitem{STP}{СТП, [Электронный ресурс]. -- Режим доступа: \url{https://www.bsuir.by/m/12_100229_1_185586.pdf}}
  \refitem{OtherRef}{Пример дополнительного источника.}
\end{ListReferences}


\end{multicols} % НЕ УДАЛЯТЬ!

% Рисунок на всю ширину страницы 
\image{
[width=15.8cm]{png/axi.png}
\caption{Пример изображения на всю ширину страницы}
}

%-%-%-%-%-%-%-%-%-%-%-%-%-%-%-%-%-%-%-%-%-%-%-%-%-%-%-%-%-%-%-%-%-%-%-%-%-%-%-%
%    __                _  _       _                                          _ 
%   /__\ _ __    __ _ | |(_) ___ | |__    _ __   ___   ___   ___   _ __   __| |
%  /_\  | '_ \  / _` || || |/ __|| '_ \  | '__| / _ \ / __| / _ \ | '__| / _` |
% //__  | | | || (_| || || |\__ \| | | | | |   |  __/| (__ | (_) || |   | (_| |
% \__/  |_| |_| \__, ||_||_||___/|_| |_| |_|    \___| \___| \___/ |_|    \__,_|
%               |___/                  
% 
%-%-%-%-%-%-%-%-%-%-%-%-%-%-%-%-%-%-%-%-%-%-%-%-%-%-%-%-%-%-%-%-%-%-%-%-%-%-%-%
\udkEN{004.8}

\topic{HARDWARE IMPLEMENTATION OF A COMPACT NEURAL NETWORK}

\informationEN{Krivaltsevich~E.~A.}

\supervisor{Vashkevich~M.~I. -- Doctor of Science}

\annotationEN{This article presents a sample template for submissions to the BSUIR Scientific Conference for Postgraduate, Master's, and Undergraduate Students. Examples of inserting images, tables, and formulas are provided, as well as recommendations for formatting references and bibliography.}

\keywordEN{SNTK, conference, template.}

\end{document}
